\chapter[一股秘溪潺潺流淌]{一股秘溪潺潺流淌}

\chaptersubtitle{《似乎在梦里见过，那样……》}

\chapternote{%
    ``Life cannot be destroyed for good, neither can history be brought entirely to a halt.
    \textbf{A secret streamlet trickles} on beneath the heavy lid of inertia and pseudo-events, slowly and inconspicuously undercutting it.
    It may be a long process, but one day it must happen: the lid will no longer hold and will start to crack.
    This is the moment when something once more begins visibly to happen, something truly new and unique...
    something truly historical, in the sense that history again demands to be heard.''
    \par
    {\raggedleft\noindent ---Václav Havel, \emph{Letter to Dr. Gustáv Husák}\par}
    \par
    “生命永远不会彻底毁灭，历史也无法完全停滞。
    在惰性与伪事件的沉重覆盖之下，\textbf{一股秘溪潺潺流淌}，缓慢而悄无声息地从下方侵蚀着这一切。
    这个过程或许漫长，但终有一日，覆盖层再难维系，开始崩裂。
    此时，某种事情再次开始显现，某种真正崭新而独特的事情...
    某种真正具有历史性的事情，因为历史再次要求人们倾听它的声音。
    \par
    {\raggedleft\noindent ——瓦茨拉夫·哈维尔《致胡萨克的公开信》\par}
    \par
    译者按：在第一集那层明亮、可爱、一切如常的表象之下，梦境、纸幕与异质空间已悄悄渗流。那便是“秘溪”最早显形的裂缝。
}

《魔法少女小圆》可理解为三个彼此分明的篇章，每个篇章聚焦于特定角色或角色组合，并探讨一个特定主题。
不过，三个主题中的每一个也都能在另外两个篇章中找到。
第一个篇章涵盖前三集，重点着墨于麻美。
透过她，作品揭示了一个根本的内在冲突：
即本剧在魔法少女类型中的定位，与其渴望实现创新突破、成为全新叙事形态之间的矛盾。
这种矛盾可进一步理解为两部相互竞争的剧集之间的斗争：
一部是本剧最初看似呈现出的、相当典型（虽略显黑暗）且深受《魔卡少女樱》影响的剧集；
另一部则是它后来演变成的、令人深感不安的《少女革命》与《新世纪福音战士》的后继之作。
简而言之，这是“虚假剧集”（以蓝光版前两集甜美温柔的片尾曲《明天再见\footnote{%
    译注：原标题为『また\, あした』，蓝光版第 1、2 话片尾曲、小圆角色歌。
}》为标志）与“真实剧集”（以第三集及之后的片尾曲 \emph{Magia} \footnote{%
    译注：``magia''源自希腊语的 ``\textgreek{μαγεία}''，意为“魔法”。
    蓝光版第 3 $\sim$ 8、11 话片尾曲，第 1、2、10 话插入曲。
}为标志）之间的对抗。
换言之，第一个篇章先在两者之间构筑一种二元对立：
一端是不熟悉编剧虚渊玄的观众大概会期待的常规魔法少女剧；
另一端是他大概会创作的黑暗的解构性剧集。
随后，这一篇章看似选择了后者。

动画伊始，也即《魔圆》本身，以定格动画这一非传统形式呈现的纸幕，缓缓升起。
艺术与非艺术之别，关键在于“框架\footnote{%
    译注：此处的“框架”兼有“将某物框定为故事或艺术”和实体“画框”之意。
}”：
一篇没有“框架”的故事只是谎言，一幅没有“框架”的画作只是涂鸦，正是“框架”定义了艺术。
这本质上是一种制度主义定义：
艺术是由艺术家创作并呈现给“艺术公众”的人工制品。
所谓“艺术公众”，是指具备艺术理论与历史知识、能识别艺术品的群体\cite{adajian-nodate-definition-art}。
换言之，艺术是能凭借自身编码方式被观众解码，并识别为艺术的存在。
由此观之，《魔圆》对定格纸片动画的运用，以及非常规动画范畴的艺术风格，或许在某种意义上正定义了《魔圆》：
它们强调这部作品自身的独特个性，而非仅仅作为某一类型的实例。

同时，纸幕令人联想到日本的“纸芝居\footnote{%
    译注：一种以连续画片配合口头讲述的日本街头表演艺术。
}”。
这种 20 世纪初的艺术形式类似于欧洲的流动木偶戏，说书人游走于城镇之间，借助便携式舞台上的连续画片讲述故事，并对漫画与动画的发展产生了深远影响\cite{nash2009raising-curtain}。
通过纸艺的运用，《魔圆》既宣示了自身对其媒介与类型源流的承继，也宣示了自身作为独立于其类型之外的独特存在。

这种诠释颇为契合，但让我们先回到剧情。
开篇的视觉母题是小圆在一个扭曲而广袤的内部空间奔跑，其棋盘格局令人联想到埃舍尔\footnote{%
    译注：莫里茨·科内利斯·埃舍尔（1898 年 6 月 17 日—1972 年 3 月 27 日），荷兰著名版画艺术家。
}笔下的错位空间。
黑白在此被严格区分，形成清晰的二元结构。
注意，这个空间并非现实。
与小圆梦境后半部分不同，它并不直接对应我们在任何时间线上看到的事件，而仅仅预示着她即将遭遇首个魔女的场所。

梦境后半部分（正如我们在动画后期所了解到的那样）实则是前一个时间线的记忆。
真正的片尾曲\emph{Magia} 在此处作为背景音乐响起，象征着这才是故事的真相。
这首曲子不会在本集或下一集的片尾播放：
电视播出版将职员表叠印在本集最后一个场景上，蓝光版则使用一首看似欢快（若不细究歌词）的歌曲，并配以欢乐的画面。
直到第三集，\emph{Magia} 才作为片尾曲登场，那一刻，作品彻底撕下了“正常”魔法少女动画的伪装。

视觉上，这段梦境以多层次灰阶构筑基调，几乎摒弃了纯黑、纯白与鲜亮色彩。
废墟都市中心，一棵巨大的枯树耸立，象征着自然与文明的双重衰亡。
瓦尔普吉斯之夜现身，她那小丑般的魔女形态与巨大齿轮暗示着人造秩序，却被癫狂的尖笑声所戳穿。
万物皆被倒置。
在《魔圆》的世界里，二元对立只是幻象，一切二元对立都坍缩为一体。
后续我们将再次探讨此点。

本集动画竭力抑制这一片段所暗示的真相。
它把这一幕轻描淡写地处理成一场梦，其手法与《魔卡少女樱》——
影响《魔圆》最明显的作品之一——
对首个预示性场景的表现如出一辙
% \footnote{%
%     《魔卡少女樱》第一集：“小樱与不可思议的魔法书”。
% }
\cite{cardcaptor-sakura-episode1-2014}。
随后的片头曲亦符合传统魔法少女动画风格（再次强调，只要忽略歌词），其中小圆的变身场景极具《美少女战士》
% \footnote{%
%     《美少女战士》第一集：“爱哭鬼小兔的华丽变身”。
% }
\cite{sailor-moon-episode1-2014}或《甜心战士》
% \footnote{%
%     《甜心战士》第一集：“一只黑色爪子抓住了心脏”。
% }
\cite{cutie-honey-episode1-2013}的特色，服装设计也令人不禁联想到小樱的几套魔法服饰。

事实上，本集其余大部分都在竭力把小圆的生活与个性渲染得尽可能无害而平凡。
她的家庭近乎完美：
一位温暖友善、在家从事园艺的父亲，一位干劲十足的职业女性母亲，还有一个惹人喜爱的小弟弟。
之后，正如《美少女战士》第一集那样，小圆叼着面包跑出家门，险些上课迟到，此类例子不胜枚举。
（实际上，这种“叼着面包赶路的迟到女生形象”在《新世纪福音战士》最后一集真嗣的幻想场景中，被用以传达“普通”绫波丽的日常生活
% \footnote{%
%     《新世纪福音战士》第二十六集：“在世界中心呼唤爱的野兽”。
% }
\cite{evangelion-episode26-2005}。）

虽然焰的出现对小圆而言打破了常态，但对观众来说却完全在预料之中。
“神秘转校生”早已是陈旧的动画母题，甚至在《凉宫春日》系列小说中被直接点名\cite{tanigawa2009haruhi}；
这些小说在一定程度上戏仿了若干动画类型，而它们都有一个共同点：
普普通通的日本青少年遭遇超自然事件。
因此，焰紧随小圆的梦境登场，又很快在高中生活的各个方面表现出众，这样的安排几乎成了陈词滥调。
她那高冷的外表\footnote{%
    译注：作者用 persona（表面形象）一词，原指戏剧面具，现指个体在不同情境中展现的公开形象或角色伪装。
}（作为一种防御性的矫饰，它兼具该词的古典与现代意义），以及她毫不费力、甚至显得百无聊赖地取得的种种优异表现，只会让她显得更加符合这种神秘设定。

除了一处例外，魔法少女世界侵入小圆生活的每一段情节都呈现出双重意义上的定式化——既缺乏新意，又符合该类型的常规套路。
她救助受伤的丘比便是又一例证，该情节与《美少女战士》中月野兔救助露娜的场景
% \footnote{%
%     《美少女战士》第一话：“爱哭鬼小兔的华丽变身”。
% }
\cite{sailor-moon-episode1-2014}高度相似。
露娜同样以释放月野兔的潜能为交换，要求她与邪恶战斗。
那么，除了我们暂且避而不谈的那一次入侵，本集中最令人耳目一新、最不落窠臼的时刻，反而是日常生活闯入其间：
小圆向仁美和沙耶香讲述自己的梦。
在本剧最不动声色却又极其精妙的场景之一中，三人认真探讨了多种对梦境的解释，从“或许小圆之前见过焰”的合理推测，到荒谬的“前世记忆”（当然，此解释最接近真相）。
无论在电视还是电影中，沙耶香甚至都指出眼下的情形简直像动画一样——
这是一个微妙的伏笔：
在后来成为魔法少女的所有角色中，她会尽最大努力去维持动画通常描绘的、理想化的“正义战士”形象。

接下来便是本集中最不像动画的场景：
魔女袭击。
这既是魔法世界对小圆安稳生活最彻底的一次侵入，也是非动画对动画最彻底的一次侵入。
定格动画与手工纸艺风格爆发式地呈现怪诞又令人不安的景象，迫使小圆和沙耶香正视她们此前一直受到保护、未曾接触的陌生存在，看见自己所知之外，那个更广阔、更狂野的世界。
异质的艺术风格强化了她们的迷失与恐惧，同时透出一丝荒诞乃至滑稽感。
尤其是“安东尼\footnote{%
    译注：蔷薇魔女的使魔，其任务是创造庭园。
}”们——那些留着胡须的棉球，类似品客薯片商标\footnote{%
    译注：品客（Pringles）为美国薯片品牌，其商标形象为一个留有褐色大胡子和红色领结的人脸。
}的形象——既荒唐又令人不安。
此种效果在后现代作品中屡见不鲜：
将来自某一语境的元素置于另一语境中，会制造出一种冲突与不协调的效果\cite{hassan1987pluralism}，而幽默\cite{boyd2004laughter-literature}与恐怖\cite{carroll1999horror-humor}皆依赖于此类不协调感。

焰除了叹道“偏偏是这种时候……”之外，对此番魔女袭击无能为力。
她知道这是不可避免的。
作为真《魔圆》的代理人（注意，她是小圆梦境中出现的唯一魔法少女），焰为维持这部更安全、更舒适的虚假动画所作的一切努力都注定失败（我们将在第十集和《叛逆》中再次见证）。
这种保护只有焰的对立项才能做到，而且也只能维持片刻。

即麻美，登场。

她刚出场就将自身定位为焰的敌人，以威胁手段保护丘比。
此时的丘比仍是那个负责唤醒女孩魔法潜能、可爱又讨喜的吉祥物兼引导者。
只有当本作撕下伪装，它才会显露其操纵人心的恶魔本质，接过反派角色。
如我们所见，麻美与焰有一个共同点：
二者皆为经验丰富且使用枪械的魔法少女。
但除此之外，她们几乎完全对立。
焰是新来的转校生，而麻美身为高年级学生，进入这所学校的时间比小圆和沙耶香更早。
焰使用的是全然普通的现代武器，麻美使用的则是魔法燧发枪。
焰的形象以直线、深色与紫色为主；
麻美的形象则以曲线（不仅体现于身材，亦见于发型）以及白色和黄色为主。
焰封闭、神秘、看似敌对；
麻美则坦率、友好。

最关键的是，麻美拥有让那部虚假的《魔圆》复原的能力，而焰则无。
在我们甚至未见其人之时，她的第一个举动便是在魔女结界中为小圆和沙耶香划定出一个安全区域。
她能够驱散真实《魔圆》的代理人，包括魔女与焰。
在动画后期，我们将看到，在另一条时间线中，麻美得知了魔女的真相——
这是“真”与“假”《魔圆》之间的关键分歧点——
她的反应既暴力又极其高效：
第一枪便击杀了冷酷无情的魔法少女杏子，第二步又制住焰，并在自己被小圆击杀前，险些把她杀了。
麻美最终被除掉，但代价极其惨重，以至该时间线（在某种意义上可视为成片的一个废弃草案）最终不得不被放弃。
若有麻美在，那真正的《魔圆》根本无从实现。

麻美在本集中的角色定位是恢复秩序，使艺术风格回归熟悉的动画常规，治愈受伤的丘比，让它得以占据类似《美少女战士》中露娜的位置，向小圆和沙耶香提供魔法力量。
凭借积极的态度、坚定的决心和强大的实力，麻美成为历代魔法少女的有力化身。
正如我们将在这一篇章中所见，她承载的是魔法少女类型中所有的经典主题。
只要她挺身守护，虚假的剧集便不会崩塌。

那就只好让她退场了。
